\section{Introduction}

Papers~7~\cite{paper7online} and~8~\cite{paper8multi} together
exhausted the easy options for online recalibration of HygienePrio's
scorer weights. Paper~7 evaluated lag-1 rolling-history grid search
and found it recovers the offline-peek ceiling at $K \in \{50, 100\}$
but \textit{harms} performance at $K = 200$. Paper~8 evaluated
multi-window smoothers (EWMA-3, trailing-mean-3) and found they
amplify the harm rather than reversing it. Paper~8 explicitly
identified adaptive procedures --- change-point detection, Bayesian
forgetting, capacity-conditional gating --- as the remaining
candidate class.

This paper evaluates the simplest adaptive procedure in that class:
a one-threshold magnitude test on the calibration-target P@50
shift between consecutive past windows. When the magnitude exceeds
$\tau = 0.05$, the procedure switches to Paper~4 fixed weights;
otherwise it uses lag-1 weights. The procedure is auditable,
trivially deployable, and contains no tuned parameters except the
pre-registered threshold.

\subsection{Contributions}

\begin{itemize}
\item \textbf{C1 --- A pre-registered evaluation of adaptive
switching.} The first pre-registered comparison of a change-point
detector against Paper~7's lag-1 baseline, Paper~8's smoothers, a
static capacity-conditional gate, and the offline-peek ceiling.

\item \textbf{C2 --- High-capacity hazard avoided.} Adaptive
switching is the first procedure in the sequence that does not harm
at $K = 200$: cell-mean delta vs \textsf{fixed} is $+0.7$~pp, and
no individual window shows a loss exceeding the $-0.01$ pre-registered
tolerance. H1 is supported.

\item \textbf{C3 --- Moderate-capacity penalty.} At $K = 50$ the
detector fires in $60\%$ of post-W1 windows and the cell-mean adaptive
P@50 is $5.3$~pp below lag-1's. The procedure trades moderate-capacity
recovery for high-capacity safety; H2 is rejected.

\item \textbf{C4 --- Marginal advantage over static gating.} Adaptive
beats a static $(K \leq 100 \to \textsf{lag1}, K \geq 200 \to
\textsf{fixed})$ gate by $+0.7$~pp at $K = 200$. H4 is supported but
the margin is small; in deployment regimes with stable known
capacity, the simpler gate may suffice.

\item \textbf{C5 --- Non-monotone firing.} Spearman$(K, \text{fire
fraction})$ is $+0.5$, well below the pre-registered $+0.8$
threshold. The detector responds to per-window distributional shifts
that are not cleanly capacity-indexed; H3 is rejected. This is itself
a useful negative result: ``fire more at higher capacity'' is the
naive expectation, but the actual firing pattern reflects
window-specific calibration-target dynamics that more sophisticated
detectors might capture.
\end{itemize}

\subsection{Relationship to Prior Papers}

This paper is the tenth in the VulnPrio research sequence. It
synthesises Paper~7's hazard finding, Paper~8's smoother rejection,
and Paper~9's~\cite{paper9selftraj} demonstration that observed
collapse patterns can be selection-induced rather than intrinsic.
The cumulative arc now reads: HygienePrio's per-pair advantage is
robust (Papers~4--6); its absolute floor is regime-dependent (Paper~6);
deployable recalibration is capacity-conditional (Paper~7); smoothing
does not help (Paper~8); Paper~6's headline collapse was a
selection-policy artifact (Paper~9); and minimal adaptive switching
finally avoids the high-capacity hazard (this paper), at the cost
of a moderate-capacity penalty and with only marginal improvement
over a static capacity gate.
